% generated by paper/make_tables.py - do not edit
\paragraph{P1\_minimal}
\begin{Verbatim}[breaklines,fontsize=\small]
Score the monetary policy stance of the following FOMC statement on a scale from -1 to +1, where -1 is dovish and +1 is hawkish.

Reply with only the number.

STATEMENT:
{statement}
\end{Verbatim}
\paragraph{P2\_defined}
\begin{Verbatim}[breaklines,fontsize=\small]
Below is a statement from the Federal Open Market Committee.

A hawkish stance leans toward tighter policy and higher interest rates, typically motivated by concern about inflation. A dovish stance leans toward looser policy and lower interest rates, typically motivated by concern about growth and employment.

Rate the stance on a continuous scale from -1 (maximally dovish) to +1 (maximally hawkish). Reply with only the number.

STATEMENT:
{statement}
\end{Verbatim}
\paragraph{P3\_analyst}
\begin{Verbatim}[breaklines,fontsize=\small]
You are a fixed income analyst. Read the FOMC statement below and tell me how hawkish it is, as a single number between -1 (dovish) and +1 (hawkish).

Respond with the number and nothing else.

STATEMENT:
{statement}
\end{Verbatim}
\paragraph{P4\_reversed} (answer sign flipped back in code)
\begin{Verbatim}[breaklines,fontsize=\small]
Assess the policy stance expressed in this FOMC statement.

Use a scale from -1 to +1 where +1 means strongly DOVISH and -1 means strongly HAWKISH.

Output only the numeric value.

STATEMENT:
{statement}
\end{Verbatim}
\paragraph{P5\_reason\_first}
\begin{Verbatim}[breaklines,fontsize=\small]
Read the FOMC statement below. Briefly note the main hawkish elements and the main dovish elements. Then give an overall stance score from -1 (dovish) to +1 (hawkish).

End your reply with a line in exactly this format:
SCORE: <number>

STATEMENT:
{statement}
\end{Verbatim}
